
\section{Related Work}

\subsection{Architectural Paradigms in Tumor Segmentation 
}For medical image segmentation, recent research shows a difference between generalized foundation models and specialized convolutional architectures. The effectiveness of an Attention-Enhanced U-Net is demonstrated by recent studies which use attention gates to suppress unimportant background features and concentrate on intricate tumor structures \cite{Islam2026}. Conversely, the SAMed-2 framework \cite{Yan2025} explores the adaptation of the "Segment Anything" architecture into a medical foundation model, leveraging temporal adapters and memory mechanisms for zero-shot generalization across datasets. Our proposed pipeline makes use of a ResUNet architecture, as opposed to the high computational requirements of foundation models like SAMed-2. This method maintains multi-scale spatial details and guarantees stable gradient flow by incorporating residual blocks into the U-Net structure. This choice validates that for specialized medical applications, a residual-based encoder-decoder provides a highly pragmatic compute-accuracy trade-off without the extensive pretraining or hardware requirements of foundation models. 

\subsection{Enhancing Classification through Localization and Transfer Learning }
Although there are many different approaches to feature extraction, deep transfer learning is still a fundamental component of tumor classification. The usefulness of applying pre-trained networks, such as Xception and DenseNet, directly on complete MRI slices is highlighted \cite{Asif2022}, which optimizes hyperparameters to attain high accuracy on benchmarks. Similar to this, an improved classification framework based on a Deep Multiple Fusion Network (DMFN) is also proposed \cite{Vure2025}. This method reduces inter-class overlap by breaking down the multi-class classification problem into a number of pairwise binary classifications using ResNet18. Rather than using complex ensemble fusions or processing the entire MRI scan, our proposed system takes a segmentation-guided multi-stage approach. The pipeline explicitly crops out irrelevant anatomical context by converting the predicted segmentation mask into a localized ROI before passing the image to a fine-tuned ResNet50 classifier. While full-image models may occasionally capture peripheral cues, this explicit localization requires the classifier to base its decision solely on the targeted tumor region. 

\subsection{Explainability, Reliability, and Summary Generation}
The transition from predictive accuracy to clinical trustworthiness depends heavily on explainable AI and reliable summary-generation mechanisms. Islam et al.~\cite{Islam2026} use Gradient-weighted Class Activation Mapping (Grad-CAM) to generate heatmaps that highlight spatial features contributing to a classification label. This method is useful for checking whether a network focuses on tumor pathology or irrelevant background structures. However, incorporating these maps into a multi-stage workflow remains challenging. Our proposed system addresses this by computing the intersection-over-union (IoU) between the segmentation mask and the Grad-CAM activation as a quantitative stage-consistency metric. This checks whether the localization and classification stages attend to the same pathology and adds an objective layer of verification to qualitative heatmaps.

Beyond visual evidence, the rise of VLMs in medical summary generation has introduced concerns about hallucinations and unconstrained generative output. The CALM-VLM framework~\cite{Dhinagar2026} addresses this issue through selective prediction: by applying temperature scaling to the model's probability distributions, it estimates uncertainty and abstains from producing a summary when prediction entropy is too high. Our proposed pipeline investigates a structural reliability alternative. Instead of relying only on statistical abstention, our VLM produces structured summaries with findings, impressions, and recommendations, while a deterministic class-conditioned fallback is used when the generated text is malformed or inconsistent with the classifier prediction. This keeps the textual explanation grounded in the verified visual prediction.

 
 